\section{Number of Ways to Arrive at Destination}
\label{sec:graph-number-of-ways-arrive}

\problemstatement{Given $n$ cities and a list of bidirectional roads
with travel times, count the number of \textbf{distinct shortest-time}
ways to travel from city $0$ to city $n{-}1$, modulo $10^9+7$.}

\begin{patternbox}
\emph{``Count the number of ways to achieve the shortest distance''}
pairs Section~\ref{sec:graph-dijkstra}'s \textit{dist} array with a
parallel \textit{ways} array, updated by the identical
reset-on-improvement, accumulate-on-tie rule the DP on LIS chapter's
Number of LIS used for counting longest chains -- here applied to
Dijkstra's relaxation instead of an array-indexed DP recurrence.
\end{patternbox}

\subsection*{The counting rule, adapted to relaxation}

When relaxing edge $u \to v$ with weight $w$, alongside the usual
distance check:
\begin{itemize}
  \item If $\textit{dist}[u] + w < \textit{dist}[v]$: a
        \textbf{strictly shorter} path to $v$ was just found through
        $u$. Update $\textit{dist}[v]$, and \textbf{reset}
        $\textit{ways}[v] \gets \textit{ways}[u]$ (every previously
        counted way to $v$ used a now-suboptimal distance).
  \item If $\textit{dist}[u] + w = \textit{dist}[v]$: this is
        \textbf{another equally short} path to $v$. \textbf{Accumulate}:
        $\textit{ways}[v] \mathrel{+}= \textit{ways}[u]$.
\end{itemize}

\subsection*{C++ implementation}

\begin{lstlisting}
#include <bits/stdc++.h>
using namespace std;
const int MOD = 1e9 + 7;

int countPaths(int n, vector<vector<int>>& roads) {
    vector<vector<pair<long long,long long>>> adj(n); // {neighbor, weight}
    for (auto& r : roads) {
        adj[r[0]].push_back({r[1], r[2]});
        adj[r[1]].push_back({r[0], r[2]});
    }

    vector<long long> dist(n, LLONG_MAX), ways(n, 0);
    dist[0] = 0;
    ways[0] = 1;
    priority_queue<pair<long long,int>, vector<pair<long long,int>>, greater<>> pq;
    pq.push({0, 0});

    while (!pq.empty()) {
        auto [d, node] = pq.top();
        pq.pop();
        if (d > dist[node]) continue;

        for (auto& [neighbor, weight] : adj[node]) {
            if (dist[node] + weight < dist[neighbor]) {
                dist[neighbor] = dist[node] + weight;
                ways[neighbor] = ways[node];
                pq.push({dist[neighbor], neighbor});
            } else if (dist[node] + weight == dist[neighbor]) {
                ways[neighbor] = (ways[neighbor] + ways[node]) % MOD;
            }
        }
    }
    return (int)(ways[n - 1] % MOD);
}

int main() {
    vector<vector<int>> roads = {{0,6,7},{0,1,2},{1,2,3},{1,3,3},
                                  {6,3,3},{3,5,1},{6,5,1},{2,5,1},{0,4,5},{4,6,2}};
    cout << countPaths(7, roads) << endl; // 4
    return 0;
}
\end{lstlisting}

\subsection*{Dry run}

\dryrun{Take a small graph: $0{-}1$ (weight $1$), $0{-}2$ (weight
$1$), $1{-}3$ (weight $1$), $2{-}3$ (weight $1$) -- two equally short
paths from $0$ to $3$.}

$\textit{dist}=[0,\infty,\infty,\infty]$, $\textit{ways}=[1,0,0,0]$.
Pop $0$: relax $1$ ($\infty\to1$, reset $\textit{ways}[1]=1$); relax
$2$ ($\infty\to1$, reset $\textit{ways}[2]=1$). Pop $1$ (dist $1$):
relax $3$ ($\infty\to2$, reset $\textit{ways}[3]=\textit{ways}[1]=1$).
Pop $2$ (dist $1$): relax $3$ -- $\textit{dist}[2]+1=2=\textit{dist}[3]$,
\textbf{tie}, accumulate $\textit{ways}[3] \mathrel{+}=
\textit{ways}[2]=1$, giving $\textit{ways}[3]=2$. Pop $3$: no further
relaxation. Final: $2$ shortest paths to node $3$ (via $1$, and via
$2$), correctly counted.

\begin{complexitybox}
\textbf{Time Complexity.} $O((V+E)\log V)$, identical to plain
Dijkstra -- the \textit{ways} bookkeeping adds only $O(1)$ work per
relaxation.
\textbf{Space Complexity.} $O(V+E)$ for the adjacency list, plus
$O(V)$ for \textit{dist} and \textit{ways}.
\end{complexitybox}

\begin{takeawaybox}
The ``exists / count / optimize'' framing introduced in the DP on
Subsequences chapter, and reused for counting longest chains in the DP
on LIS chapter, extends naturally beyond dynamic programming tables --
any algorithm computing an optimum via relaxation or comparison (here,
Dijkstra's distance updates) can be paired with a parallel counting
array using the exact same reset-on-improvement,
accumulate-on-tie rule.
\end{takeawaybox}
